% Chapter 15: Tenses in Writing

\chapter{Tenses in Writing: Keeping Time Steady}

\section{Pedagogical Purpose}
This chapter builds on the understanding of basic tenses by teaching students how to use them consistently in a piece of writing. It introduces the present continuous tense and focuses on avoiding unnecessary tense shifts, which is a common problem for young writers.

\begin{learningobjectives}
The learner can:
\begin{itemize}
    \item identify the tense used in a passage and explain why it fits the purpose.
    \item write a short passage describing a daily routine using present tense.
    \item write a short narrative about the past using consistent past tense.
    \item use "will" and "going to" to talk about plans and predictions.
    \item form and use the present continuous tense correctly.
    \item find and correct an unnecessary tense shift in a passage.
\end{itemize}
\end{learningobjectives}

\section{Prerequisite Check}
\begin{quickcheck}
State the tense of each sentence.
\begin{enumerate}
    \item Rohan \textbf{plays} cricket every Saturday.
    \item Maya \textbf{visited} her cousin last month.
    \item We \textbf{will go} to the fair next week.
\end{enumerate}
\end{quickcheck}

\begin{rememberbox}
In Chapter 14 you learned three time zones: present, past, and future. Now you will learn how to use them properly across a whole piece of writing — and meet one new form, the present continuous.
\end{rememberbox}

\section{Chapter Opener}
Rohan has written about his day for the class newsletter, but Ms. Sharma notices something odd.

\textbf{Rohan's Newsletter Entry:} "I wake up at seven. Yesterday I brush my teeth and I will eat breakfast. Then I am going to school and I played with my friends."

\textbf{Ms. Sharma:} "Rohan, this is lively writing, but your times are jumping all over the place! One moment you're in the past, then the future, then the present. Let's fix your time-travelling paragraph."

\textbf{Maya:} "It's like Leo's diary muddle from before, but in a whole paragraph!"

\section{Language Experience}
\textbf{Two Versions of Maya's Weekend}

> \textbf{Version A (Muddled):} On Saturday, I woke up late and I eat breakfast. My mother is calling me and I go downstairs. We watched a movie and now we are having lunch.

> \textbf{Version B (Steady):} On Saturday, I woke up late and I ate breakfast. My mother called me and I went downstairs. We watched a movie and then we had lunch.

\section{Look and Notice}
\begin{thinkbox}
\begin{enumerate}
    \item Both versions describe the same weekend. Which one feels easier to follow?
    \item In Version A, how many different tenses can you find?
    \item In Version B, which single tense is used all the way through?
    \item Why might switching tense confuse a reader?
\end{enumerate}
\end{thinkbox}
\section{Core Explanation}
When we write more than one sentence, we usually keep the \textbf{same tense} throughout, unless the meaning truly changes in time. This is called \textbf{tense consistency}.

\begin{itemize}
    \item Writing about a \textbf{routine or fact} \(\rightarrow\) stay in the \textbf{present tense}.
    \item Telling a \textbf{story about the past} \(\rightarrow\) stay in the \textbf{past tense}.
    \item Describing \textbf{plans or predictions} \(\rightarrow\) use \textbf{"will"} or \textbf{"going to"}.
    \item Describing something happening \textbf{right now} \(\rightarrow\) use the \textbf{present continuous} ("is/are" + verb-\textbf{ing}).
\end{itemize}

\subsection*{Present Continuous Tense}
The present continuous tense is used for actions happening right now. It is formed with \textbf{am/is/are} + the base verb + \textbf{-ing}.

\begin{examplebox}
    \begin{itemize}
        \item I \textbf{am reading}.
        \item She \textbf{is writing}.
        \item They \textbf{are playing}.
    \end{itemize}
\end{examplebox}

\begin{thinkbox}
    What is the difference between "I play football" and "I am playing football"?
    
    "I play football" is a general statement or habit (Simple Present).
    "I am playing football" means the action is happening right now (Present Continuous).
\end{thinkbox}

\section{Worked Example}
\begin{examplebox}
\textbf{Passage to fix:} "Every day, Maya feeds her fish. Yesterday, she cleans the tank and she will change the water."

\textbf{Let's Think Step-by-Step:}
\begin{enumerate}
    \item \textbf{Check the second sentence.} "Yesterday" signals the past, but the verb "cleans" is present tense — this is an error.
    \item \textbf{Check the third sentence.} "She will change" is future, but there is no future time signal — likely meant to continue the past story.
\end{enumerate}

\textbf{Answer:} "Every day, Maya feeds her fish. Yesterday, she \textbf{cleaned} the tank and she \textbf{changed} the water."

\textbf{Why:} The first sentence stays present because it describes a routine; the second and third sentences must both be past because "yesterday" sets the whole story in the past.
\end{examplebox}

\section{Guided Practice}
\begin{activitybox}{Activity A: Identify the Tense}
Read each passage and name the main tense used.
\begin{enumerate}
    \item "Rohan gets up at six. He brushes his teeth. He eats breakfast." (Present - routine)
    \item "Last Sunday, we visited the zoo. We saw elephants. We had ice cream." (Past - narrative)
    \item "Look! The children are playing in the park. They are laughing loudly." (Present Continuous - happening now)
\end{enumerate}
\end{activitybox}

\begin{activitybox}{Activity B: Completion}
Fill in the correct present continuous form.
\begin{enumerate}
    \item Look! The dog \_\_\_ (bark) at the gate.
    \item Right now, Maya \_\_\_ (write) a letter.
\end{enumerate}
\end{activitybox}

\begin{activitybox}{Activity C: Correct the Tense}
Rewrite the following passage with the correct and consistent tense.
"Yesterday was a holiday. We go to the beach. The sun shines brightly and we swim in the sea. Later, we will build a big sandcastle."
\end{activitybox}

\section{Independent Practice}
\begin{activitybox}{Activity B: Correct the Tense}
Rewrite the following passage with the correct and consistent tense.

"Yesterday was a holiday. We go to the beach. The sun shines brightly and we swim in the sea. Later, we will build a big sandcastle."
\end{activitybox}

\section{Summary}
\begin{summarybox}
    \begin{itemize}
        \item Keep the \textbf{same tense} throughout a passage unless the time truly changes.
        \item Use \textbf{present tense} for routines and facts.
        \item Use \textbf{past tense} for stories that have already happened.
        \item Use \textbf{"will"} or \textbf{"going to"} for plans and predictions.
        \item Use the \textbf{present continuous} ("is/am/are" + verb-\textbf{ing}) for actions happening right now.
    \end{itemize}
\end{summarybox}